% Gerado de tcc.md (exportado do Claude Docs) por ferramentas/gera_latex.php em 12/11/2026 21:45.
% Para mudar o texto, edite este arquivo; se gerar de novo a partir do documento, as edições feitas aqui se perdem.

\noindent The Individualized Education Plan (PEI, in Portuguese) sets learning goals for students eligible for special education, but tracking those goals usually relies on notes written by the teacher, with no data showing progress over time. This work presents UNIA (from the Portuguese Unir Necessidades Individuais e Aprendizagem, uniting individual needs and learning), a web system that helps teachers track the PEI goals of children in the early school years and measures their performance through mini games. The proposal builds on AppPEI, a digital educational resource built as a spreadsheet at the Rio de Janeiro State University, from which the four areas of the skills matrix and the three support levels were adopted: does not perform, performs with help, and performs without support. The research is applied and exploratory, combining a literature review with incremental software development in PHP, following the Model-View-Controller (MVC) pattern, with a MySQL database, tested on the MariaDB 10.4 server bundled with XAMPP. In UNIA, the teacher registers the student and their goals, and the student plays five-round sessions in 12 mini game formats, spread across 24 activities, six per area. The system checks each answer and records correctness, hint use and response time; at the end, it classifies the session's support level and computes, for each goal, the mean accuracy, the share of sessions without support, the share of sessions with hints and the change between the first and second halves of the sessions, shown in tables and accessible charts. Technical verification included 63 automated tests of the rules, all passing, and an end-to-end test on a temporary database; a heuristic evaluation based on Nielsen’s principles found ten usability problems, two of them major, which were then fixed. The system has not yet been evaluated with teachers and students; the classification rule and the parameters adopted are prototype hypotheses, and their pedagogical validation is the next step. The results indicate that automatic performance measurement can be integrated into PEI monitoring, turning one-off records into series that show the progress of each goal.

\vspace{\baselineskip}
\noindent\textbf{Keywords:} Individualized Education Plan. Special education. Digital games. Learning monitoring. Web system.
